\section*{Chapter \ref{ch:rc:commodity-and-energy-derivatives} --- Commodity and Energy Derivatives}

\begin{solution}{exo:rc:commodity-and-energy-derivatives:1}
27.2\%, against 68.7\% at delivery. The short-term factor's loading $e^{-\kappa\tau}$ has
decayed to $e^{-1.5} = 0.22$: a forward a year out moves mainly with the long-term level,
because shocks to today's supply and demand are expected to fade before it delivers.
\end{solution}

\begin{solution}{exo:rc:commodity-and-energy-derivatives:2}
Every schedule buys and sells at the same forward price and pays the injection and
withdrawal costs, so the best deterministic plan is to do nothing. The facility is still
a portfolio of \omterm{def:rc:commodity-and-energy-derivatives:spread}{spread options}: as prices move, spreads between months open and the owner
can exploit them; its value is all extrinsic.
\end{solution}

\begin{solution}{exo:rc:commodity-and-energy-derivatives:3}
Its intrinsic value, discounted one day: $3.50-2.74-0.50 = 0.26$ dollars per MMBtu
(0.2600 after a day's discounting at 4\%).
\end{solution}

\begin{solution}{exo:rc:commodity-and-energy-derivatives:4}
The spread's volatility is $\sqrt{\sigma_1^2+\sigma_2^2-2\rho\sigma_1\sigma_2}$ (at zero strike): higher correlation
means a less volatile spread and a cheaper option. Close delivery months are highly
correlated (0.970 for July and January at July's expiry in the chapter), distant ones less.
\end{solution}

\begin{solution}{exo:rc:commodity-and-energy-derivatives:5}
It starts with the intrinsic plan hedged with forwards, and changes the plan only when the
new plan, valued on the new curve, is worth more than the old one after re-hedging:
each change adds non-negative value locked in by forwards.
\end{solution}

\begin{solution}{exo:rc:commodity-and-energy-derivatives:6}
The strip gives the right to exercise every monthly call; the swing caps the total, so its
holder must choose which months to use, giving up some exercises the strip would take.
The strip is an upper bound (6.11 against 4.37 in the chapter).
\end{solution}

\begin{solution}{exo:rc:commodity-and-energy-derivatives:7}
At zero strike Kirk's formula is Margrabe's, exact for lognormal forwards: the difference
from Monte Carlo is 0.007 cent, within the noise. At 0.50 it is $-0.375$ cent: the sum
$F_2+K$ is not lognormal, and the approximation of its volatility by $\sigma_2F_2/(F_2+K)$ is
off near the money, where the option is most sensitive to the distribution's shape.
\end{solution}

\begin{solution}{exo:rc:commodity-and-energy-derivatives:8}
A book hedged at intrinsic is hedged against parallel moves of the forward curve, but it
owns options: in February 2021 the right to withdraw fast at spot was worth millions of
dollars to a fast facility, and a book that did not use it gave it away. Conversely, a
book short flexibility (supply contracts with swing) had a large loss. The intrinsic hedge
does not measure either.
\end{solution}

\begin{solution}{pb:rc:commodity-and-energy-derivatives:1}
\textbf{1.} USD~634\,447, 699\,550 and 698\,182.
\textbf{2.} About USD~64\,000, 10\% of the intrinsic: the right to change the schedule as
spreads between months move.
\textbf{3.} The short-term volatility $\sigma_s$, with the mean reversion $\kappa$: they set how much
near months move against each other.
\textbf{4.} It lets the facility act once a month on monthly prices; a salt cavern acts daily,
on daily prices whose spikes average out in a monthly price.
\textbf{5.} Buy the summer forwards and sell the winter forwards of the intrinsic plan, in
the plan's volumes.
\textbf{6.} 2.88 dollars on 1 February; 23.86 on 17 February.
\textbf{7.} 2.62 dollars.
\textbf{8.} 500\,000 MMBtu on five days (11, 12, 16, 17 and 18 February), a gain of
USD~4\,306\,870.
\textbf{9.} 90\,000 MMBtu on nine days, a gain of USD~479\,057.
\textbf{10.} The cavern runs out of gas after five days and sells only on the best ones; the
reservoir sells a little every day, including the days of 3.35 to 3.76 dollars.
\textbf{11.} $\ln(23.86/3.76) = 1.85$ against $0.687\sqrt{4/252} = 0.087$: 21 standard deviations,
impossible in the model.
\textbf{12.} Jumps (a jump-diffusion or regime-switching spot) or a stochastic volatility
that rises with scarcity.
\textbf{13.} The spike lasted a week and left the forward curve for the next season almost
unchanged, so the seasonal spread (the intrinsic value) did not move; the options gained
because they were exercised in the spike.
\textbf{14.} Swing contracts and other rights to take gas at a fixed price, firm pipeline
capacity out of the frozen region, and call options on daily gas.
\textbf{15.} Sellers of those rights: suppliers with swing obligations who had to buy gas
at spot to deliver at the contract price, and generators short fuel for their power sales.
\textbf{16.} The value sits in rare, extreme weeks; the average year understates it, and a
single event can pay for years of fees.
\textbf{17.} The intrinsic position and its hedges, the \omterm{def:rc:commodity-and-energy-derivatives:extrinsic}{extrinsic value} and its sensitivity to
volatility, the value under historical stress weeks, and the physical limits (rates,
ratchets, contracts that could oblige deliveries).
\textbf{18.} Add spikes (jumps or regimes) to the short-term factor, and decide daily.
\textbf{19.} Named result: \emph{the salt cavern in February 2021}: the fast facility gains
USD~4\,306\,870 on 500\,000 MMBtu, the slow one USD~479\,057 on 90\,000.
\textbf{20.} The options to move gas through time at short notice, which pay most in
exactly the weeks no model of the average predicts.
\end{solution}

\begin{solution}{iq:rc:commodity-and-energy-derivatives:1}
Short-dated forwards are more volatile than long-dated ones. In a two-factor model a
mean-reverting short-term factor moves near contracts with loading $e^{-\kappa\tau}$, fading
for distant deliveries, while a long-term factor moves them all.
\iqlookfor{the observation and the mechanism.}
\end{solution}

\begin{solution}{iq:rc:commodity-and-energy-derivatives:2}
Zero strike: Margrabe, exact for lognormal prices. Non-zero strike: \omterm{def:rc:commodity-and-energy-derivatives:kirk}{Kirk's approximation},
the Bjerksund--Stensland refinement, one-dimensional integration conditioning on one price,
or Monte Carlo. Kirk is accurate for small strikes relative to the second price and less so
near the money or for large strikes and very different volatilities.
\iqlookfor{Margrabe, Kirk and when to integrate.}
\end{solution}

\begin{solution}{iq:rc:commodity-and-energy-derivatives:3}
Intrinsic: optimise the schedule on today's forwards and hedge it. \omterm{def:rc:commodity-and-energy-derivatives:extrinsic}{Rolling intrinsic}:
re-optimise and re-hedge as the curve moves, a lower bound on full optionality. Full: a
stochastic dynamic programme on spot (least-squares Monte Carlo or a tree), valuing daily
decisions under constraints.
\iqlookfor{the three layers and the ordering.}
\end{solution}

\begin{solution}{iq:rc:commodity-and-energy-derivatives:4}
Simulate spot; backward induction on an inventory grid; per period and end inventory,
regress realised future value on basis functions of the state; act on the estimates, carry
realised values. Biases: in-sample foresight (high), suboptimal regression policies (low);
use independent paths for the valuation pass and check against \omterm{def:rc:commodity-and-energy-derivatives:extrinsic}{rolling intrinsic}.
\iqlookfor{the algorithm and the two biases.}
\end{solution}

\begin{solution}{iq:rc:commodity-and-energy-derivatives:5}
Hedge the intrinsic with forwards and re-hedge as the plan changes; measure \omterm{def:rc:commodity-and-energy-derivatives:extrinsic}{extrinsic
value} and its sensitivity to volatility; stress with historical spikes and cold weeks; track
physical constraints and inventory; limit the short flexibility sold in swing contracts;
keep liquidity for margin on the hedges.
\iqlookfor{hedging, stress and physical limits.}
\end{solution}

\begin{solution}{iq:rc:commodity-and-energy-derivatives:6}
Gaussian factors give lognormal prices with thin tails; spikes come from inelastic supply
and demand and last days. Add jumps with fast mean reversion, regime switching, or a
fundamental model of supply and demand with a steep stack.
\iqlookfor{thin tails and a remedy with fast reversion.}
\end{solution}
